
This work sits at the intersection of meta-learning, benchmark-driven evaluation, and automated machine learning. Prior research assumes dataset characteristics are available for meta-learning, but this work identifies metadata extraction as a practical bottleneck.

\textbf{Meta-Learning and Algorithm Selection.} Meta-learning approaches learn to optimize learning algorithms across tasks. AutoML systems like Auto-sklearn and TPOT use dataset meta-features (statistical properties) to warm-start Bayesian optimization or evolutionary search. These systems assume dataset characteristics are readily computable, extracting features like sample count, dimensionality, class balance, and skewness from in-memory data. Rice formalized algorithm selection as mapping problem instances to optimal algorithms, introducing the Algorithm Selection Problem framework. Smith-Miles developed instance space analysis to visualize problem hardness and algorithm performance regions. These frameworks assume features are known, focusing on predictive modeling rather than the practical challenge of feature extraction from heterogeneous sources.

This work investigates whether published benchmark results can train a meta-classifier without requiring dataset downloads. This scenario is relevant for practitioners evaluating method choices before committing computational resources, or researchers synthesizing guidance from existing literature. Prior meta-learning work does not address the extraction bottleneck: how to collect dataset characteristics when datasets are distributed across repositories, APIs, and paper supplements.

\textbf{Benchmark Studies.} Recent benchmark papers provide empirical evidence of method performance variability but do not aggregate findings into predictive models. Zhou et al. evaluate nine medical imaging datasets in federated learning settings, reporting that no single algorithm (FedAvg, FedProx, SCAFFOLD, FedDyn) consistently achieves top performance. Small datasets (TB with 668 samples) benefit from DDPM+LS augmentation (+17 percentage points) while large datasets (ColonPath with 10,000 samples) show minimal improvement (+0.3 percentage points). These results suggest dataset size correlates with augmentation effectiveness, but the correlation is not tested explicitly.

Champneys et al. establish baseline comparisons for nonlinear system identification tasks, finding that NARX-Poly achieves 0.032 RMSE on Wing-Hing saturation versus 0.126 for LSTM. Structured, low-dimensional physics-informed problems favor polynomial bases over recurrent architectures. Other domain-specific benchmarks (OGB for graphs, FedML for federated learning, LEAF for federated settings) document method rankings on specific datasets but do not extract generalizable patterns. Benchmark papers provide ground truth for meta-learning (method rankings across datasets) but do not attempt to learn predictors from this data, and metadata required for prediction is absent from published results.

\textbf{Method Selection Heuristics.} Practitioners rely on informal heuristics for method selection. Afkanpour et al. conduct a systematic review of federated learning challenges, concluding that ''data heterogeneity matters'' and practitioners should ''consider data structure.'' Liao et al. characterize heterogeneity in federated learning, providing taxonomies of non-IID data types (label skew, feature skew, quantity skew) but no decision rules for matching data types to methods. These guidelines are qualitative and unactionable---''consider data structure'' does not specify which features to measure or how to map features to method families. Domain folklore ($vision\rightarrow CNN$, time-$series\rightarrow RNN, tabular\rightarrow tree$ ensembles) persists despite evidence that exceptions are common.

\textbf{Positioning.} This work's contribution is procedural rather than algorithmic. It does not propose a novel meta-learning architecture or feature extraction method. Instead, it identifies and quantifies a practical bottleneck: literature-based metadata extraction yields only 14\% average coverage for critical dataset characteristics. This finding explains why prior meta-learning work focuses on AutoML scenarios (where datasets are already loaded) rather than literature synthesis scenarios (where datasets must be collected). The contribution is exposing a hidden infrastructure assumption in meta-learning research and proposing two-stage data collection (identification + characteristic extraction) as a field-wide need.
